\section{Evaluation framework and triage method}\label{sec:method}

\subsection{Problem setting}
A flow-based NIDS observes a stream of flow records $x\in\mathbb{R}^d$ and must decide which flows to surface to a security operations centre (SOC). Each flow has a binary label $y\in\{0,1\}$ (benign/attack) and, for attacks, a family label $f$. A detector $h$ trained on a \emph{source} dataset outputs a score $s(x)=\hat{P}(y=1\mid x)$. At deployment the detector meets a \emph{target} stream whose distribution can differ from the source in two ways that we study separately: (i)~\emph{cross-dataset shift}, where the network, capture tooling and attack implementations change (benign and attack distributions both move); and (ii)~\emph{unseen-family shift}, where the network is the same but an attack family absent from training appears. We call the setting with neither shift \emph{in-domain}.

We take the SOC's point of view and assume that analysts can inspect at most a fraction $b$ of the flows in a stream (the \emph{review budget}), and that attacks are rare. Unless stated otherwise the target streams have attack prevalence $\pi=1\%$ and budgets are $b\in\{1\%,2\%,5\%\}$; a perfect ranker therefore reaches full recall for every $b\ge\pi$.

\subsection{Datasets and harmonisation}\label{sec:data}
We use five public benchmarks (Table~\ref{tab:data}). Three of them, CIC-IDS2017~\cite{sharafaldin2018toward}, CSE-CIC-IDS2018 (the AWS-hosted ``Processed Traffic Data for ML Algorithms'' release) and ACI-IoT-2023, were produced with CICFlowMeter and share 76 flow features plus the transport protocol; we align them column by column (the 2018 naming is canonical) and verify the alignment programmatically with a token-level name check. These three datasets support cross-dataset experiments in both directions. UNSW-NB15~\cite{moustafa2015unsw} (official training and testing partitions, merged) and the network part of ToN-IoT have their own feature sets and are used for in-domain and unseen-family experiments. We remove identifiers that act as shortcuts or do not transfer (flow IDs, IP addresses, ports, timestamps), one-hot encode low-cardinality categorical fields (UNSW-NB15 and ToN-IoT, top-15 values per field), replace non-finite values by zero, and remove exact duplicates over features and label. To keep the study computationally balanced, we then cap benign flows at 200{,}000 and each attack family at 30{,}000 per dataset by uniform sampling; for CSE-CIC-IDS2018, benign flows are first Bernoulli-subsampled at 5\% while streaming the 16.2~million rows. Families are also mapped to coarse groups (DoS, DDoS, Scan, BruteForce, WebAttack, Botnet, Infiltration) for the cross-dataset analysis. The exact provenance (URLs, commit hashes, checksums) is listed in the supplementary repository.

The CIC-IDS2017 copy we could obtain is the stratified subsample released by the mlp4nids project (all attack flows, roughly a quarter of the benign flows), and the ToN-IoT copy retains 50{,}000 of the 300{,}000 normal records of the official train--test file. Because we resample to fixed prevalences for every operational metric, these subsampling choices affect the benign diversity available for training, not the metrics' definitions; we return to this in Section~\ref{sec:threats}.

\subsection{Detectors and calibration}
We evaluate four widely used model families: LightGBM~\cite{ke2017lightgbm} (300 trees, 63 leaves), random forest~\cite{breiman2001random} (150 trees), a multilayer perceptron (128--64 units, early stopping) and logistic regression, all from scikit-learn~\cite{pedregosa2011scikit} or the LightGBM library. Inputs are transformed with a signed $\log(1+|x|)$ followed by standardisation fitted on the training split. Each source dataset is split, stratified by family, into training (50\%), calibration (20\%) and test (30\%) partitions; all results are repeated over five seeds that change both the split and the model initialisation.
Probabilities are post-processed with Platt scaling~\cite{platt1999probabilistic} or isotonic regression~\cite{zadrozny2002transforming} fitted on the source calibration split, and we report expected calibration error (ECE, 15 equal-width bins)~\cite{naeini2015obtaining,guo2017calibration} and the Brier score.

\subsection{Operating points and base rates}
Following common practice, the alert threshold $\tau_\alpha$ is set on the source calibration split so that the benign false-alarm rate is $\alpha=1\%$. On a target stream we measure the realised true-positive rate (TPR) and false-positive rate (FPR) at $\tau_\alpha$, the \emph{FPR inflation} $\mathrm{FPR}/\alpha$, and the precision the detector would achieve at attack prevalence $\pi$,
\begin{equation}
\mathrm{Prec}_\pi=\frac{\mathrm{TPR}\,\pi}{\mathrm{TPR}\,\pi+\mathrm{FPR}\,(1-\pi)},
\end{equation}
which makes the base-rate effect~\cite{axelsson2000base} explicit.

\subsection{Budgeted conformal triage}\label{sec:triage}
Ranking flows by $s(x)$ exploits what the detector has learned but cannot surface traffic that resembles nothing seen in training, which is precisely where unseen attacks and shifted traffic live~\cite{sommer2010outside}. We therefore formulate triage as testing, for each flow, the null hypothesis
\[
H_0(x):\ \text{``$x$ is benign and typical of the source traffic''},
\]
and review the flows with the strongest evidence against $H_0$. Two conformal p-values~\cite{vovk2022algorithmic,angelopoulos2023conformal,bates2023testing} are computed against the \emph{benign} flows of the source calibration split, $\mathcal{C}_0=\{x_i\}_{i=1}^{n}$:
\begin{align}
p_{\mathrm{det}}(x) &= \frac{1+\#\{i: s(x_i)\ge s(x)\}}{n+1}, &
p_{\mathrm{nov}}(x) &= \frac{1+\#\{i: u(x_i)\ge u(x)\}}{n+1},
\end{align}
where $u(x)$ is a novelty score, here the mean Euclidean distance to the $k=5$ nearest neighbours among up to 20{,}000 \emph{benign} training flows in the transformed feature space. $p_{\mathrm{det}}$ is small when the detector recognises attack-like behaviour; $p_{\mathrm{nov}}$ is small when the flow is atypical of the benign traffic the model was trained on, whether or not it resembles a known attack. Using benign flows as the reference matches the null hypothesis; in an ablation it also separated held-out families from benign traffic better than a reference that included training attacks (Section~\ref{sec:discussion}). We combine them either with a Bonferroni rule, $p_{\mathrm{B}}=\min\{1,2\min(p_{\mathrm{det}},p_{\mathrm{nov}})\}$, or with Fisher's method, $p_{\mathrm{F}}=1-F_{\chi^2_4}\!\big(-2\log p_{\mathrm{det}}-2\log p_{\mathrm{nov}}\big)$. If benign target flows are exchangeable with $\mathcal{C}_0$, each $p_{\mathrm{det}}$ and $p_{\mathrm{nov}}$ is super-uniform, and the Bonferroni rule $p_{\mathrm{B}}\le\alpha$ flags at most a fraction $\alpha$ of benign flows in expectation, whatever the dependence between the two scores. Fisher's rule is exact only under independence, so we use it purely as a ranking. Cross-dataset shift breaks exchangeability by design; one purpose of our experiments is to measure how much of the guarantee survives.

\paragraph{Budget enforcement without target labels} Given a budget $b$, the review queue consists of the flows whose priority $-\log p(x)$ exceeds a threshold $t_b$. We compare two ways of choosing $t_b$: (a)~\emph{source transfer}, the $(1-b)$-quantile of priorities on a source stream with the same prevalence; and (b)~\emph{label-free warm-up}, the $(1-b)$-quantile on the first 20\% of the unlabeled target stream, applied to the remaining 80\%. Option (b) needs no labels because a budget is a statement about volume, not about accuracy. We report the realised \emph{load} (queue size divided by $b\cdot N$; 1 means the budget is respected exactly) and the attack recall inside the queue. We also report the recall of the exact top-$b$ fraction of each stream, which isolates ranking quality from threshold selection.

\subsection{Experimental protocols}
\paragraph{Protocol XD (in-domain and cross-dataset)} For every source dataset, seed and model, we train on the source training split, calibrate on the source calibration split, and evaluate on the source test split (in-domain) and, for the three CICFlowMeter datasets, on the test split of each other dataset (six directed cross-dataset pairs). No target information is used except the unlabeled warm-up window of the budget rule.
\paragraph{Protocol OS (unseen families)} For every dataset and every attack family with at least 300 flows, we remove that family from the training and calibration splits and evaluate on the test split. The target stream contains all benign test flows plus the held-out family at 1\% prevalence, emulating a new campaign. We report the recall of the held-out family at $\tau_\alpha$, the fraction of its flows that the isotonic-calibrated detector labels benign with confidence above 95\% (\emph{confident misses}), and the triage metrics. This protocol uses LightGBM and three seeds.
